\section{INTRODUCTION}
In this paper, we study a class of nonsmooth nonconvex constrained
optimization problems formulated as
\begin{equation}
\tag{{\bf P}}
\label{prob:main}
\begin{alignedat}{3}
   \min\limits_{\bx \in \RR^n} \quad
      & F(\bx) &&:= \phi(\bx) - g(\bx),
%       \\
%    \text{s.t.} \quad
%       & \bc_{\calE}(\bx) &&= \zero, \\
%       & \bc_{\calI}(\bx) &&\leq \zero ,
\end{alignedat}
\end{equation}
where the objective $F = \phi - g$ has a \emph{difference-of-convex}
(DC) structure: $\phi:\RR^n\to\RR\cup\{+\infty\}$ and
$g:\RR^n\to\RR\cup\{+\infty\}$ are proper closed convex functions, and
$g$ is allowed to be nonsmooth, so that $F$ is in general both
nonsmooth and nonconvex. 
% The constraint mappings
% $\bc_{\calE}:\RR^n\to\RR^{m_{\calE}}$ and
% $\bc_{\calI}:\RR^n\to\RR^{m_{\calI}}$ are continuously differentiable
% and define a nonempty feasible set
% $\calX = \{\bx \in \RR^n : \bc_{\calE}(\bx) = \zero,\
% \bc_{\calI}(\bx) \le \zero\}$.

Two settings motivate our study of Problem~\eqref{prob:main}. The first is deterministic constrained DC programming, in which both
$\phi$ and $g$ can be evaluated exactly. DC functions form a vector
space and constitute a remarkably rich subclass of locally Lipschitz
nonconvex functions \citep{Zhang-2020-Complexity,kong2024cost,davis2022gradient}; the formulation subsumes smooth nonconvex minimization ($g \equiv 0$) as a
trivial instance and covers a broad spectrum of classical
applications, including nonconvex quadratic programming
\citep{lethi1997solving,lethi1998branch}, convex--concave programming
\citep{yuille2003concave,shen2016disciplined}, kernel selection \citep{argyriou2006dc},
and bilevel programming \citep{lethi2009dc}. While the unconstrained DC
objective is by now well understood, the nonlinear constraints
$\bc_{\calE}$ and $\bc_{\calI}$, which arise naturally from
structural or resource requirements such as energy budgets in deep
network compression \citep{yang2019ecc,chen2018constraint} and safety
constraints in control \citep{paternain2019constrained},introduce
severe geometric complications of their own.

A second motivation for \eqref{prob:main} is regularized empirical risk minimization in large-scale machine learning. The DC structure of the objective naturally captures the combination of a data-fitting risk and a nonconvex regularizer. Specifically, any smooth nonconvex loss can be reformulated as a DC function by adding and subtracting a strongly convex quadratic term. Moreover, popular nonconvex regularizers such as the capped-$\ell_1$ penalty \citep{zhang2010analysis} and SCAD \citep{fan2001variable} admit explicit DC decompositions $r = r_{\mathrm{cvx}} - g$. In large-scale training, the data-fitting component in $\phi$ is a
population risk accessible only through stochastic gradients under
data sampling, whereas the subtracted component $g$, which collects
the concave part of the regularizer, remains a fixed deterministic
function. 

To ground our framework, we provide two concrete motivating examples that can be formulated as instances of Problem~\eqref{prob:main}:

\subsection{Motivating Examples}

\begin{enumerate}
  \item \textbf{Deterministic DC objective.}
Losses given by a pointwise minimum of convex functions are
intrinsically DC, via the identity
$\min\limits_{k\le K} a_k=\sum_{k=1}^{K}a_k-\max\limits_{k\le K}\sum_{j\ne k}a_j$
for convex $a_1,\dots,a_K$. For instance, in piecewise-affine
(clusterwise) regression, each observation is explained by the best of
$K$ candidate affine models,
\[
   \min\limits_{\bw_1,\dots,\bw_K}\ \sum_{i=1}^{N}\,
   \min\limits_{1\le k\le K}\bigl|y_i-\langle \bw_k,\bx_i\rangle\bigr|,
\]
and applying the identity samplewise yields $F=\phi-g$ with
$\phi(\bw)=\sum_{i,k}|y_i-\langle \bw_k,\bx_i\rangle|$ and
$g(\bw)=\sum_{i}\max\limits_{k}\sum_{j\ne k}|y_i-\langle \bw_j,\bx_i\rangle|$.
Both components are convex, piecewise affine, and nonsmooth, with
subgradients available in closed form.

  \item \textbf{Stochastic nonconvex regularized ERM.}
  Consider a convex data-fitting loss $\ell$ coupled with a capped-$\ell_1$ penalty:
  \[
     \min\limits_{\bx}\ \underbrace{\EE_{\xi}[\ell(\bx,\xi)] + \lambda\|\bx\|_1}_{\phi(\bx)}
     \;-\;\underbrace{\lambda\sum_i\max\limits\{|x_i|-\theta,0\}}_{g(\bx)}.
  \]
  Here, the data-fitting component $\phi$ is a population risk evaluated solely via stochastic gradients, whereas the subtracted penalty $g$ is a fixed, deterministic, and nonsmooth convex function whose subgradients can be computed in closed form.
\end{enumerate}

\subsection{Related Work}
\label{subsec:known-and-open}

\paragraph{Non-asymptotic rates for DC optimization.}

Removing differentiability of the subtracted component $g$ has been a
central theme in DC programming. Recent non-asymptotic analyses have
made this theme quantitative, but their guarantees are expressed through
different decomposition-dependent certificates. For DCA and CCCP-type
methods, rates are usually stated in terms of objective decrease, step
lengths, or critical-distance proxies under additional smoothness or
curvature assumptions
\citep{yurtsever2022cccp,abbaszadehpeivasti2023rate}. Projection-free
DC methods over compact convex domains instead use Frank--Wolfe-type
linear minimization gaps \citep{khamaru2018convergence,maskan2025}.
Stochastic DC methods use critical-distance criteria or one-sided gap
functions, often with a smooth or variance-reduced stochastic component
\citep{xu2019stochastic,le2022stochastic,nguyen2025improved}. These
certificates are value-typed or algorithm-dependent, while the exact DC
residual
\[
    \operatorname{dist}\bigl(\zero,\partial\phi(\bx)-\partial g(\bx)\bigr)
\]
is norm-typed and same-point. Passing from a gap bound to a residual
bound therefore changes units, typically through a square-root relation.
Existing finite-time results hence do not settle whether exact
same-point $\epsilon$-criticality is itself attainable for nonsmooth DC
pairs. The first question is therefore:

\emph{Is $\epsilon$-criticality finite-time attainable in DC
optimization? If not, does a unified single notion exist that relaxes it while
remaining both finite-time attainable and sound?}

\paragraph{Non-asymptotic rates for unconstrained nonsmooth nonconvex
optimization.}
For general Lipschitz objectives, $\epsilon$-Clarke stationarity is
unattainable in finite time \citep{Zhang-2020-Complexity}, while the
structure-blind Goldstein $(\delta,\epsilon)$-stationarity relaxation admits
$O(\delta^{-1}\epsilon^{-3})$ guarantees
\citep{Zhang-2020-Complexity,davis2022gradient},
one $\epsilon$-factor away from the $O(\epsilon^{-2})$ smooth
benchmark. These results treat the objective as a black box. 

The second question is therefore:

\emph{Under an finite-sum attainable targat of DC
criticality, can the DC structure be provably exploited to match the
$O(\epsilon^{-2})$ rate of unconstrained smooth nonconvex
optimization?}

\paragraph{Constrained nonconvex optimization.}
The rates above are established either without constraints or over
simple convex compact sets accessed through projections or linear
minimization oracles \citep{khamaru2018convergence,maskan2025}; general nonlinear
constraints are handled by penalty or augmented-Lagrangian schemes
whose inner solves make the methods double-loop and couple the inner
accuracy to the outer complexity
\citep{xie2021complexity,lu2022penalty,yu2021convergence,kanzow2025adaptive}. The third question is:

\emph{can a single-loop method, free of inner solvers and
backtracking, retain the optimal $O(\epsilon^{-3})$ complexity under general
nonlinear, possibly nonconvex, constraints? (ALM)}

\subsection{Contributions}
We tackle the three questions for Problem~\eqref{prob:main} when both components of the DC decomposition may be nonsmooth.  This class is rich enough to cover
constrained DC programs and nonconvex-regularized empirical risk
minimization. At first sight, the DC structure makes the search for
stationary points easier on both sides of the oracle--target ledger:
the natural target, DC criticality
$\zero\in\partial\phi(\bx)-\partial g(\bx)$, is strictly weaker than
Clarke stationarity, and the natural oracle, returning values and
subgradients of $\phi$ and $g$ separately, is strictly stronger than
a Clarke subgradient oracle. Surprisingly, even with a weaker target
and a stronger oracle, finding an $\epsilon$-critical point, i.e., a
point $\bar\bx$ with
$\operatorname{dist}(\zero,\partial\phi(\bar\bx)-\partial g(\bar\bx))
\le\epsilon$, is intractable. In other words, no algorithm
interacting with the component-wise DC oracle can guarantee an
$\epsilon$-critical point among its queried iterates within a finite
number of oracle calls.This intractability suggests that, to obtain meaningful
non-asymptotic results for DC problems, we need to refine the notion
of criticality. We then introduce such a unified notion and base our analysis on
it, leading to the following main contributions:

\begin{itemize}
\item We show that a traditional $\epsilon$-critical point cannot be
obtained in finite time, even under the strong component-wise DC
oracle, returning values and subgradients of $\phi$ and $g$
separately.


\item We study $(\delta_\phi,\delta_g,\epsilon)$-critical points
(\Cref{def:delta-eps-critical}), which unify existing relaxations of
DC criticality.For smooth DC functions, this notion reduces to the usual
$\epsilon$-stationarity by choosing \(\delta_\phi=O(\epsilon^2/L_\phi)\) and \(\delta_g=O(\epsilon^2/L_g)\).

\item For unconstrained DC optimization, we show that proximal DCA returns a \((0,\delta_g,\epsilon)\)-critical point within
\(\mathcal O(\max\limits\{\epsilon^{-2},\delta_g^{-1}\})\) subproblem solves.
When \(g\) is additionally \(L_g\)-smooth, one may take
\(\delta_g=\mathcal O(\epsilon^2/L_g)\), so that the two-sided certificate
collapses to the standard \(\mathcal O(\epsilon^{-2})\) guarantee for
\(\epsilon\)-stationarity. \blue{Moreover, composing this certificate with an
inner solver of per-subproblem cost \(T_{\rm sub}(\omega)\) at accuracy
\[
\omega=
\mathcal O\!\left(
\min\limits\left\{
\eta\epsilon^2,\,
\frac{\delta_\phi^2}{\eta G_\phi^2},\,
\frac{\delta_g^2}{\eta G_g^2}
\right\}
\right)
\]
yields a \((\delta_\phi,\delta_g,\mathcal O(\epsilon))\)-critical point with total
inner-iteration complexity
\[
\mathcal O\!\left(
\max\limits\left\{
\frac{\Delta}{\eta\epsilon^2},\,
\frac{\eta\Delta G_g^2}{\delta_g^2}
\right\}
\cdot T_{\rm sub}(\omega)
\right),
\]
by the same inexact bookkeeping as in the bare-DCA case
(\Cref{cor:dca-inexact-target}); here \(T_{\rm sub}(\omega)\) is the number of
inner iterations to solve one proximal subproblem to accuracy \(\omega\).}

\item We extend the criticality notion to the constrained problem
\eqref{prob:main} and propose a single-loop method attaining a
$(\delta_\phi,\delta_g,\epsilon)$-critical point in $O(\epsilon^{-2})$
iterations under general nonlinear, possibly nonconvex, constraints
(see \Cref{sec:complexity-analysis-DCA}), matching the unconstrained rate above
despite nonsmoothness in both components.
\end{itemize}


\subsection{Notations}
Throughout, $\|\cdot\|$ denotes the Euclidean norm and
$\langle\cdot,\cdot\rangle$ the standard inner product on $\RR^n$. For
a proper closed convex function $h$, $\partial h(\bx)$ denotes its
convex subdifferential, $\partial_\delta h(\bx)$ its
$\delta$-subdifferential, and $h^*$ its Fenchel conjugate; for a
locally Lipschitz $F$, $\partial_C F(\bx)$ denotes the Clarke
subdifferential. We write
$\operatorname{dist}(A,B) := \inf\limits_{\ba\in A,\,\bb\in B}\|\ba-\bb\|$
for nonempty sets $A,B\subseteq\RR^n$, and $\operatorname{dist}(\bv,A)
:= \operatorname{dist}(\{\bv\},A)$. For two sets \(A,B\subseteq\RR^n\), we use the Minkowski difference
\(A-B\), also called the pairwise difference set, defined by
\[
    A-B:=\{\ba-\bb:\ba\in A,\ \bb\in B\}.
\]  Bold lowercase letters denote
vectors and $\zero$ the zero vector. 

